% ============================================================
%  Rozdział 8 — Regulamin Szkolenia
% ============================================================
\section{Regulamin Szkolenia z Podstaw Kajakarstwa Górskiego 2026}

% ------------------------------------------------------------
\subsection*{§1. Organizator}
\addcontentsline{toc}{subsection}{§1. Organizator}
% ------------------------------------------------------------

Organizatorem szkolenia jest \textbf{Wrocławski Klub Kajakowy „PrzeWrotka"} (dalej:
Organizator).

% ------------------------------------------------------------
\subsection*{§2. Charakter szkolenia}
\addcontentsline{toc}{subsection}{§2. Charakter szkolenia}
% ------------------------------------------------------------

\begin{itemize}[leftmargin=1.5em, itemsep=2pt]
  \item Szkolenie ma charakter sportowy i obejmuje zajęcia teoretyczne oraz praktyczne,
    w tym zajęcia basenowe oraz spływy kajakowe.
  \item Ze względu na intensywność ćwiczeń wymagających dobrej kondycji fizycznej oraz
    stabilnego stanu psychicznego, uczestnicy zobowiązani są do przestrzegania niniejszego
    regulaminu.
\end{itemize}

% ------------------------------------------------------------
\subsection*{§3. Ochrona danych osobowych (RODO)}
\addcontentsline{toc}{subsection}{§3. Ochrona danych osobowych (RODO)}
% ------------------------------------------------------------

\begin{itemize}[leftmargin=1.5em, itemsep=2pt]
  \item Dane osobowe uczestników są przetwarzane przez Organizatora wyłącznie w celu
    realizacji szkolenia oraz zapewnienia bezpieczeństwa uczestników, zgodnie z RODO
    (Rozporządzenie PE i Rady UE 2016/679).
  \item Podanie danych osobowych jest dobrowolne, lecz niezbędne do udziału w szkoleniu.
  \item Uczestnik ma prawo dostępu do swoich danych, ich poprawiania oraz żądania usunięcia
    zgodnie z obowiązującymi przepisami prawa.
\end{itemize}

% ------------------------------------------------------------
\subsection*{§4. Wizerunek}
\addcontentsline{toc}{subsection}{§4. Wizerunek}
% ------------------------------------------------------------

Uczestnik wyraża zgodę na nieodpłatne utrwalanie i wykorzystywanie swojego wizerunku
na zdjęciach i nagraniach wykonanych podczas szkolenia przez Organizatora, kadrę lub
współuczestników, w celach promocyjnych i informacyjnych (strona internetowa, media
społecznościowe), zgodnie z obowiązującymi przepisami prawa.

% ------------------------------------------------------------
\subsection*{§5. Opłaty}
\addcontentsline{toc}{subsection}{§5. Opłaty}
% ------------------------------------------------------------

Koszt szkolenia \textbf{obejmuje}:

\begin{itemize}[leftmargin=1.5em, itemsep=2pt]
  \item udział w zajęciach teoretycznych i praktycznych (basen, spływy),
  \item opiekę instruktorów,
  \item korzystanie z klubowego sprzętu kajakowego (kajak, wiosło, fartuch, kamizelka
    asekuracyjna, kask, ewentualnie odzież specjalistyczna),
  \item ubezpieczenie NNW na czas trwania kursu.
\end{itemize}

Koszt szkolenia \textbf{nie obejmuje}:

\begin{itemize}[leftmargin=1.5em, itemsep=2pt]
  \item dojazdu na zajęcia i spływy,
  \item kosztu noclegu,
  \item odzieży i sprzętu biwakowego.
\end{itemize}

Wpłata pierwszej raty jest równoznaczna z akceptacją niniejszego regulaminu. W przypadku
rezygnacji po rozpoczęciu szkolenia (tj.\ od dnia pierwszego wykładu) pierwsza rata
nie podlega zwrotowi. Szczególne przypadki rozpatrywane są indywidualnie przez Zarząd
WKK „PrzeWrotka".

% ------------------------------------------------------------
\subsection*{§6. Odpowiedzialność uczestnika}
\addcontentsline{toc}{subsection}{§6. Odpowiedzialność uczestnika}
% ------------------------------------------------------------

\begin{itemize}[leftmargin=1.5em, itemsep=2pt]
  \item Uczestnik bierze udział w szkoleniu na własną odpowiedzialność, mając świadomość
    swojego stanu zdrowia oraz ewentualnych przeciwwskazań medycznych.
  \item Każde pogorszenie samopoczucia lub stan zdrowia mogący wpływać na bezpieczeństwo
    należy niezwłocznie zgłosić instruktorowi.
\end{itemize}

% ------------------------------------------------------------
\subsection*{§7. Zasady bezpieczeństwa}
\addcontentsline{toc}{subsection}{§7. Zasady bezpieczeństwa}
% ------------------------------------------------------------

\begin{itemize}[leftmargin=1.5em, itemsep=2pt]
  \item Podczas zajęć obowiązuje bezwzględny zakaz spożywania alkoholu, środków
    odurzających i psychotropowych oraz przebywania pod ich wpływem.
  \item Uczestnicy zobowiązani są do wykonywania poleceń instruktorów. Ewentualne uwagi
    lub pytania należy zgłaszać w bezpiecznym momencie, po wykonaniu polecenia.
  \item Podczas zajęć na wodzie uczestnik podlega bezpośrednio instruktorowi prowadzącemu
    grupę.
  \item Pływanie poza zajęciami z instruktorem wymaga jego zgody i odbywa się wyłącznie
    na wskazanym akwenie z odpowiednią asekuracją.
\end{itemize}

% ------------------------------------------------------------
\subsection*{§8. Warunki uczestnictwa}
\addcontentsline{toc}{subsection}{§8. Warunki uczestnictwa}
% ------------------------------------------------------------

\begin{itemize}[leftmargin=1.5em, itemsep=2pt]
  \item Warunkiem udziału w spływach jest dostarczenie podpisanego regulaminu przed
    rozpoczęciem zajęć terenowych.
  \item Organizator zastrzega sobie prawo do wykluczenia uczestnika ze szkolenia
    w przypadku:
    \begin{itemize}[leftmargin=1.5em, itemsep=1pt]
      \item nieprzestrzegania niniejszego regulaminu,
      \item stwierdzenia spożycia alkoholu lub środków odurzających,
      \item stwarzania zagrożenia dla siebie lub innych uczestników,
      \item innych istotnych powodów wskazanych przez Organizatora.
    \end{itemize}
\end{itemize}

% ------------------------------------------------------------
\subsection*{§9. Zmiany organizacyjne}
\addcontentsline{toc}{subsection}{§9. Zmiany organizacyjne}
% ------------------------------------------------------------

Kierownictwo spływów ma prawo odwołać lub przerwać zajęcia (np.\ z powodu warunków
pogodowych, hydrologicznych lub bezpieczeństwa), proponując nowy termin realizacji.

% ------------------------------------------------------------
\subsection*{§10. Postanowienia końcowe}
\addcontentsline{toc}{subsection}{§10. Postanowienia końcowe}
% ------------------------------------------------------------

Prawo interpretacji niniejszego Regulaminu przysługuje Zarządowi WKK „PrzeWrotka".

\vspace{2cm}
\noindent
